\documentclass{report}
\usepackage{hyperref}
\usepackage{amsmath}
\usepackage{amsfonts}
\usepackage{bm}
\usepackage{geometry}
\usepackage{natbib}
\bibliographystyle{aasjournal}
\usepackage{graphicx}
\usepackage{tikz}
\geometry{a4paper, margin=2.5cm}

\hypersetup{
    colorlinks=true,
    citecolor=blue,
    linkcolor=black,
    pdfborder={0 0 0}
}


\makeatletter
\let\jnl@style=\rm
\def\ref@jnl#1{{\jnl@style#1}}

\def\aj{\ref@jnl{AJ}}
\def\actaa{\ref@jnl{Acta Astron.}}
\def\araa{\ref@jnl{ARA\&A}}
\def\apj{\ref@jnl{ApJ}}
\def\apjl{\ref@jnl{ApJ}}
\def\apjs{\ref@jnl{ApJS}}
\def\ao{\ref@jnl{Appl.~Opt.}}
\def\apss{\ref@jnl{Ap\&SS}}
\def\aap{\ref@jnl{A\&A}}
\def\aapr{\ref@jnl{A\&A~Rev.}}
\def\aaps{\ref@jnl{A\&AS}}
\def\azh{\ref@jnl{AZh}}
\def\baas{\ref@jnl{BAAS}}
\def\bac{\ref@jnl{Bull. astr. Inst. Czechosl.}}
\def\caa{\ref@jnl{Chinese Astron. Astrophys.}}
\def\cjaa{\ref@jnl{Chinese J. Astron. Astrophys.}}
\def\icarus{\ref@jnl{Icarus}}
\def\jcap{\ref@jnl{J. Cosmology Astropart. Phys.}}
\def\jrasc{\ref@jnl{JRASC}}
\def\memras{\ref@jnl{MmRAS}}
\def\mnras{\ref@jnl{MNRAS}}
\def\na{\ref@jnl{New A}}
\def\nar{\ref@jnl{New A Rev.}}
\def\pra{\ref@jnl{Phys.~Rev.~A}}
\def\prb{\ref@jnl{Phys.~Rev.~B}}
\def\prc{\ref@jnl{Phys.~Rev.~C}}
\def\prd{\ref@jnl{Phys.~Rev.~D}}
\def\pre{\ref@jnl{Phys.~Rev.~E}}
\def\prl{\ref@jnl{Phys.~Rev.~Lett.}}
\def\pasa{\ref@jnl{PASA}}
\def\pasp{\ref@jnl{PASP}}
\def\pasj{\ref@jnl{PASJ}}
\def\rmxaa{\ref@jnl{Rev. Mexicana Astron. Astrofis.}}
\def\qjras{\ref@jnl{QJRAS}}
\def\skytel{\ref@jnl{S\&T}}
\def\solphys{\ref@jnl{Sol.~Phys.}}
\def\sovast{\ref@jnl{Soviet~Ast.}}
\def\ssr{\ref@jnl{Space~Sci.~Rev.}}
\def\zap{\ref@jnl{ZAp}}
\def\nat{\ref@jnl{Nature}}
\def\iaucirc{\ref@jnl{IAU~Circ.}}
\def\aplett{\ref@jnl{Astrophys.~Lett.}}
\def\apspr{\ref@jnl{Astrophys.~Space~Phys.~Res.}}
\def\bain{\ref@jnl{Bull.~Astron.~Inst.~Netherlands}}
\def\fcp{\ref@jnl{Fund.~Cosmic~Phys.}}
\def\gca{\ref@jnl{Geochim.~Cosmochim.~Acta}}
\def\grl{\ref@jnl{Geophys.~Res.~Lett.}}
\def\jcp{\ref@jnl{J.~Chem.~Phys.}}
\def\jgr{\ref@jnl{J.~Geophys.~Res.}}
\def\jqsrt{\ref@jnl{J.~Quant.~Spec.~Radiat.~Transf.}}
\def\memsai{\ref@jnl{Mem.~Soc.~Astron.~Italiana}}
\def\nphysa{\ref@jnl{Nucl.~Phys.~A}}
\def\physrep{\ref@jnl{Phys.~Rep.}}
\def\physscr{\ref@jnl{Phys.~Scr}}
\def\planss{\ref@jnl{Planet.~Space~Sci.}}
\def\procspie{\ref@jnl{Proc.~SPIE}}

\let\astap=\aap
\let\apjlett=\apjl
\let\apjsupp=\apjs
\let\applopt=\ao

\title{{\it mirrordisk}:  Mirror-Symmetry Residual Analysis of Protoplanetary Disks in ALMA Spectral-Line Cubes}
\author{Masataka Aizawa (\texttt{aizw.masa@gmail.com})}
\date{Apr 28, 2026}

\begin{document}

\maketitle

\tableofcontents

\chapter{Overview}
{\it mirrordisk} is a Python package for extracting non-axisymmetric residual structure from ALMA spectral-line cubes by comparing each channel with its mirrored counterpart across the disk minor axis. In the three-dimensional space of sky coordinates $(x, y) = (\mathrm{RA}, \mathrm{Dec})$ and velocity/channel $v$, an approximately axisymmetric disk is expected to satisfy
\begin{equation}
   I(v, x, y) = I(2v_{\rm sys} - v, x_{\rm sym}, y_{\rm sym}),
\end{equation}
where $(x_{\rm sym}, y_{\rm sym})$ is the reflection of $(x, y)$ across the disk minor axis and $v_{\rm sys}$ is the systemic velocity. The minor axis is specified by the disk position angle and a center offset.

In the current implementation, $x$ and $y$ correspond to $\Delta$RA and $\Delta$Dec in arcsec, respectively. Following the usual astronomical image convention, positive $\Delta$RA is shown to the left, while positive $\Delta$Dec is shown upward. The mirror axis is represented as a straight line in the image plane,
\begin{equation}
    y = \tan(\mathrm{PA}) (x - x_{\rm cen}),
\end{equation}
where PA is the position angle and $x_{\rm cen}$ is the $x$-intercept of the minor axis, namely the $\Delta$RA offset where the line crosses $y=0$. Thus, the fitted geometric parameters are $(x_{\rm cen}, \mathrm{PA})$, rather than a full two-dimensional disk-center coordinate.
Figure~\ref{fig:mirror_symmetry} illustrates this mirror geometry, including the minor axis, the intercept $x_{\rm cen}$, and the reflected coordinate pair $(x, y) \leftrightarrow (x_{\rm sym}, y_{\rm sym})$.

\begin{figure}[t]
\centering
\begin{tikzpicture}[scale=1.1, >=stealth]
    % axes
    \draw[->] (3.6,0) -- (-3.2,0) node[left] {$x=\Delta \mathrm{RA}$};
    \draw[->] (0,-2.4) -- (0,3.0) node[above] {$y=\Delta \mathrm{Dec}$};

    % minor axis
    \draw[thick, blue] (-0.6,-1.97) -- (2.4,2.48) node[above right, black] {minor axis};

    % x_cen
    \filldraw[blue] (0.73,0) circle (1.4pt);
    \draw[dashed, blue] (0.73,0) -- (0.73,-0.45);
    \node[below, blue] at (0.73,-0.45) {$x_{\rm cen}$};

    % PA arc
    \draw[->, blue] (1.55,0) arc[start angle=0,end angle=56,radius=0.82];
    \node[blue] at (1.68,0.55) {$\mathrm{PA}$};

    % point and reflected point
    \filldraw[red] (2.7,0.4) circle (1.6pt);
    \node[right, red] at (2.7,0.4) {$(x,y)$};
    \filldraw[red] (0.36,1.98) circle (1.6pt);
    \node[left, red] at (0.36,1.98) {$(x_{\rm sym},y_{\rm sym})$};

    % connection
    \draw[dashed, red] (2.7,0.4) -- (0.36,1.98);

    % midpoint / perpendicular
    \filldraw[black] (1.53,1.19) circle (1.1pt);
    \draw[black] (1.45,1.09) -- (1.54,1.03) -- (1.60,1.12);
\end{tikzpicture}
\caption{Schematic of the mirror symmetry used by \textit{mirrordisk}. A point $(x,y)$ at velocity $v$ is compared with its reflected position $(x_{\rm sym}, y_{\rm sym})$ across the minor axis at velocity $2v_{\rm sys}-v$.}
\label{fig:mirror_symmetry}
\end{figure}

The code reconstructs the mirror cube
\begin{equation}
    I_{\rm mirror}(v, x, y) = I(2v_{\rm sys} - v, x_{\rm sym}, y_{\rm sym}),
\end{equation}
and separates the symmetric and asymmetric components as
\begin{align}
     I(v, x, y) &= I_{\rm sym}(v, x, y) + I_{\rm asym}(v, x, y), \\
I_{\rm sym}(v, x, y) &= \frac{I(v, x, y) + I_{\rm mirror}(v, x, y)}{2}, \\
I_{\rm asym}(v, x, y) &= \frac{I(v, x, y) - I_{\rm mirror}(v, x, y)}{2}.
\end{align}
In this interpretation, $I_{\rm sym}$ traces the approximately axisymmetric disk component, while $I_{\rm asym}$ isolates non-axisymmetric residual structure.

The current runnable workflow is based on \texttt{tests/main.py}, configured by \texttt{tests/config.py} and a target CSV table. The pipeline provides the following steps:
\begin{itemize}
\item \texttt{mirror}: fit mirror-symmetry parameters and generate mirrored and residual FITS cubes.
\item \texttt{channel\_maps}: create channel-by-channel PDF comparison plots.
\item \texttt{html}: build an interactive 3D residual isosurface HTML.
\item \texttt{residuals}: produce residual moment maps and higher-order scalar maps.
\end{itemize}

Common command-line options for every step are:
\begin{itemize}
\item \texttt{--targets-csv}: choose the target table.
\item \texttt{--output-root}: choose the output base directory.
\item \texttt{--dataset}: run one dataset only.
\end{itemize}

The installation, configuration, and runnable examples are summarized in the following chapters.


\chapter{Install \label{sec:reuired}}
Install the package with
\begin{verbatim}
pip install .
\end{verbatim}
or, for editable development,
\begin{verbatim}
pip install -e .
\end{verbatim}

The default pipeline requires the following Python packages:
\begin{itemize}
    \item numpy
    \item scipy
    \item astropy
    \item matplotlib
    \item plotly
    \item pandas
    \item psutil
    \item astrodendro
\end{itemize}
These dependencies are listed in \texttt{setup.py} and \texttt{requirements.txt}. JAX-backed modules are optional and are not required for the default pipeline.

\chapter{Input and Configuration \label{sec:prep_sec}}

\section{Target CSV format \label{sec:format_data}}
The pipeline expects a target CSV table describing one or more FITS spectral-line cubes. The CSV must include the following columns:
\begin{itemize}
    \item \textbf{source\_name}: dataset identifier used by \texttt{--dataset} and output naming.
    \item \textbf{source\_fits\_path}: full path to the input FITS cube.
    \item \textbf{pa\_deg\_initial}: initial position angle in degrees.
    \item \textbf{x\_cen\_initial}: initial center offset in arcsec.
    \item \textbf{vsys\_initial}: initial systemic velocity in km/s. This may be left blank to estimate it from the cube.
\end{itemize}

An example is
\begin{verbatim}
source_name,source_fits_path,pa_deg_initial,x_cen_initial,vsys_initial
J1615,/path/to/J1615_cube.fits,146,0,4.7
\end{verbatim}

The repository example is \texttt{tests/targets.csv}. The current examples assume local FITS paths; the repository does not bundle the data cubes.

\section{Configuration in \texttt{tests/config.py}}
\texttt{tests/config.py} combines target-specific CSV values with shared defaults.

\subsection{Mirror fitting}
\texttt{MIRROR\_DEFAULT\_PARAMS} controls the mirror fitting step:
\begin{itemize}
    \item The fitted parameters are \textbf{systemic velocity} $v_{\rm sys}$, \textbf{center offset} $x_{\rm cen}$, and \textbf{position angle} PA.
    \item Here, $x_{\rm cen}$ is the \textbf{$x$-intercept of the minor axis}, that is, the \textbf{$\Delta$RA offset} in arcsec where the minor-axis line crosses $y=0$. It is not a full two-dimensional disk-center coordinate.
    \item Their initial values come from the target CSV columns \texttt{vsys\_initial}, \texttt{x\_cen\_initial}, and \texttt{pa\_deg\_initial}.
    \item \textbf{inp\_max}: this parameter determines the spatial extent used in the mirror fitting. The interpolation grid is built over $-\,\texttt{inp\_max} \le x, y \le \texttt{inp\_max}$ in arcsec, so it sets the half-width of the fitted square region on the sky.
    \item \textbf{n\_points}: this parameter determines the spatial sampling of the fitting grid. It is the number of interpolation samples on each spatial axis, so larger values produce a finer grid for the fitting.
    \item \textbf{vec\_offsets}: this parameter determines the velocities used in the mirror fitting. These are velocity offsets in km/s relative to \texttt{vsys\_initial}, and the fitting velocities are computed as $v = \texttt{vsys\_initial} + \texttt{vec\_offsets}$.
\end{itemize}

\subsection{Channel maps}
\texttt{CHANNEL\_MAP\_PARAMS} controls output for channel-map plotters
\begin{itemize}
    \item \textbf{dv}: this parameter determines the velocity sampling of the channel-map output. Channel maps are requested at intervals of \texttt{dv} in km/s around the fitted systemic velocity.
    \item \textbf{vlim}: this parameter determines the total velocity range shown in the channel maps. The code samples channels over approximately $-\,\texttt{vlim} \le v-v_{\rm sys} \le \texttt{vlim}$.
    \item \textbf{plot\_lim}: this parameter determines the spatial field of view shown in each channel-map panel. The displayed axes are limited to $\pm\,\texttt{plot\_lim}$ arcsec.
    \item \textbf{font\_size}, \textbf{axes\_labelsize}: these parameters determine the annotation size in the PDF figures, including titles and axis labels.
\end{itemize}

\subsection{3D HTML residual view}
\texttt{HTML\_PARAMS} controls the residual isosurface export:
\begin{itemize}
    \item \textbf{skip1}: this parameter determines the spatial downsampling applied before building the 3D residual cube. The code keeps every \texttt{skip1}-th spatial pixel along both sky axes.
    \item \textbf{skip2}: this parameter determines the spectral downsampling applied before meshing. The code keeps every \texttt{skip2}-th velocity channel.
    \item \textbf{vlim}: this parameter determines the velocity range retained for the 3D residual view. Only channels satisfying $|v-v_{\rm sys}| < \texttt{vlim}$ are used.
    \item \textbf{plot\_lim}: this parameter determines the spatial region retained for the 3D residual view. Only pixels satisfying $|x| < \texttt{plot\_lim}$ and $|y| < \texttt{plot\_lim}$ are used.
    \item \textbf{contour\_sigma}: this parameter determines the isosurface threshold. The plotted contour level is computed from the measured noise as $\texttt{contour\_sigma} \times \sigma$.
\end{itemize}

\subsection{Residual analysis}
\texttt{RESIDUAL\_ANALYSIS\_PARAMS} controls the post-processing maps:
\begin{itemize}
    \item \textbf{skip1}: this parameter determines the spatial downsampling used before residual moment analysis. The code keeps every \texttt{skip1}-th spatial pixel.
    \item \textbf{skip2}: this parameter determines the spectral downsampling used before residual moment analysis. The code keeps every \texttt{skip2}-th velocity channel.
    \item \textbf{vlim}: this parameter determines the velocity range included in the residual analysis. Only channels satisfying $|v-v_{\rm sys}| < \texttt{vlim}$ are analyzed.
    \item \textbf{plot\_lim}: this parameter determines the spatial region included in the residual analysis. Only pixels satisfying $|x| < \texttt{plot\_lim}$ and $|y| < \texttt{plot\_lim}$ are analyzed.
    \item \textbf{min\_npix}: this parameter determines the minimum connected structure size retained by the dendrogram mask. Smaller connected residual features are discarded.
    \item \textbf{nsigma}: this parameter determines the significance threshold for the dendrogram mask. Residual emission must exceed approximately $\texttt{nsigma} \times \sigma$ to be included.
\end{itemize}

\chapter{Run scripts \label{sec:run_script}}
\section{Quick start}
\subsection{Overview}
If no step is given, \texttt{tests/main.py} runs all steps in this order:
\begin{itemize}
    \item \texttt{mirror}
    \item \texttt{channel\_maps}
    \item \texttt{html}
    \item \texttt{residuals}
\end{itemize}

\subsection{Example commands}
Run one step:
\begin{verbatim}
python tests/main.py \
  --targets-csv tests/targets.csv \
  --output-root /path/to/output \
  mirror
\end{verbatim}

Run all steps for all datasets listed in the CSV:
\begin{verbatim}
python tests/main.py \
  --targets-csv tests/targets.csv \
  --output-root /path/to/output
\end{verbatim}

Run one dataset only when needed:
\begin{verbatim}
python tests/main.py \
  --targets-csv tests/targets.csv \
  --output-root /path/to/output \
  --dataset J1615
\end{verbatim}

To run a single step, append the step name such as \texttt{mirror}, \texttt{channel\_maps}, \texttt{html}, or \texttt{residuals} at the end of the command.

The options are:
\begin{itemize}
    \item \textbf{--targets-csv}: CSV file defining input FITS cubes and initial parameters.
    \item \textbf{--output-root}: base output directory.
    \item \textbf{--dataset}: dataset name from the target CSV. If omitted, every dataset in the CSV is processed.
\end{itemize}

\subsection{Output structure}
Given \texttt{--output-root /path/to/output}, the pipeline writes:
\begin{itemize}
    \item \texttt{/path/to/output/result/mirror/}: fitted parameters, mirrored FITS, residual FITS.
    \item \texttt{/path/to/output/result/channel\_maps/<source\_name>/}: channel-by-channel PDF figures.
    \item \texttt{/path/to/output/html/}: interactive residual isosurface HTML.
    \item \texttt{/path/to/output/moms/}: moment maps and higher-order residual maps.
\end{itemize}

Mirror products are named from the input FITS stem:
\begin{itemize}
    \item \texttt{<stem>\_mirrorparams.npy}
    \item \texttt{<stem>\_mirror.fits}
    \item \texttt{<stem>\_res.fits}
\end{itemize}

\section{Step details}
\subsection{\texttt{mirror}}
\begin{itemize}
    \item Fits the mirror-symmetry parameters \textbf{$v_{\rm sys}$}, \textbf{$x_{\rm cen}$}, and \textbf{PA}.
    \item In the current implementation, $x_{\rm cen}$ means the \textbf{minor-axis $x$-intercept}, namely the $\Delta$RA position where the minor axis crosses $y=0$.
    \item The initial guesses are taken from \texttt{vsys\_initial}, \texttt{x\_cen\_initial}, and \texttt{pa\_deg\_initial} in the target CSV.
    \item Writes optimized parameters and generates mirrored and residual FITS cubes.
\end{itemize}

\subsection{\texttt{channel\_maps}}
\begin{itemize}
    \item Loads the original, mirrored, and residual cubes.
    \item Writes one PDF per sampled velocity channel with side-by-side panels.
\end{itemize}

\subsection{\texttt{html}}
\begin{itemize}
    \item Builds an interactive 3D isosurface view from the residual cube.
    \item Writes an HTML file under the HTML output directory.
\end{itemize}

\subsection{\texttt{residuals}}
\begin{itemize}
    \item Masks significant residual structure with \texttt{astrodendro}.
    \item Produces positive and negative residual moment maps.
    \item Produces additional scalar maps for moment 8 and moment 9.
\end{itemize}

\section{Notes}
\begin{itemize}
    \item The examples assume \texttt{pip install .} or \texttt{pip install -e .} has already been run, so \texttt{PYTHONPATH=src} is not needed.
    \item \texttt{--dataset} is optional. If omitted, every row in the CSV is processed.
    \item If \texttt{vsys\_initial} is blank in the CSV, the code estimates the systemic velocity from the cube.
\end{itemize}


\bibliography{ref.bib}

\end{document}
